% Conclusion and future work (owner SK). The limitations section is deliberately
% longer than the contributions section; every item in it is measured, not feared.
\section{Conclusion and Future Work}
\label{sec:conclusion}

An English-trained deepfake detector that reaches 0.85\% EER on ASVspoof does not
work on Hindi--English telephone speech. It reaches 48.21\% EER on clean code-mixed
audio and 58.42\% over the channel, with AUC below chance in the channel condition,
and its failure mode is the rejection of genuine speakers rather than the acceptance
of clones. Low-rank adaptation over a frozen encoder, training 1.13\% of the
parameters on speaker-disjoint code-mixed data and its clones, brings channel-matched
error to 2.75\% EER while retaining 2.42\% on clean English. Full fine-tuning on identical
data is better on the tool it saw (0.83\%) and collapses to chance on English; the
adapters keep both, which is what a deployment needs.

The detector runs on a live call. It scores four-second windows of the caller's
speech, escalates only when consecutive windows agree, and warns the callee --- never
the caller --- about ten seconds into the speech, using the same operating point the
evaluation reports.

\subsection{What we could not do}

\textbf{Against an unseen tool, over a telephone line, we cannot show that this system
works.} On clean audio the adapters transfer to Tortoise-TTS at 5.44--12.01\% EER
against matched floors of 29.49--31.43\%, which is learning beyond the shortcut. Over
the channel the deployed model averages 35.59\% across five training seeds against a
matched floor of 37.79\% and beats it in no seed at the speaker level; its
per-speaker error ranges from 6.4\% to 45.5\%, and the published single run
(30.99\%) was a favourable draw. Only S3 trained on XTTS alone stays below its floor
there in every seed (19.72\% against 33.64\%), and adding RVC to its training removes
that. Randomising recording cues during training brings the deployed recipe to
26.63\% over five seeds, but that variant was chosen after a test score had been seen
and still does not clear the floor. We therefore have no conclusive evidence that the
deployed model detects an unseen tool on a real phone line. This is the sharpest
limitation of the work and it constrains every deployment claim made here.

\textbf{New voices speaking freely are rejected.} Every code-mixed test set is read
lecture audio. On 24 new speakers talking spontaneously (HiACC), the deployed model
calls 54.37\% of genuine clips fake, against 0.89\% for the read lectures, and every
speaker is above 10\%. Adaptation taught the model what genuine lecture recordings
sound like, not genuine Hindi--English speech in general.

\textbf{The deployed model is for Hindi--English calls only.} English passed through
the same telephone channel rises to 24.20\% EER, against 6.50\% for the English-only
S1 (Sec.~\ref{sec:tradeoff}); the 2.42\% English figure holds for clean audio.

\textbf{The attacks are ours, and they are not hard.} We generated the code-mixed
attacks because no benchmark provides them, which means our numbers are against
attacks of our own construction, screened by our own quality checks. The team rated
the seen attack 1.5 out of 5 for sounding human. Detector numbers against easy fakes
are optimistic.

\textbf{The telephone leg is simulated, not carried.} The channel condition is an
8~kHz G.711 render with additive noise, validated by measurement and by ear, and the
live system is implemented against a real telephony carrier --- but the carrier leg
has not been exercised on live numbers. The protocol this work was designed around
is the telephone network, and that is the one leg not yet demonstrated end to end.

\textbf{One live-call behaviour is unexplained.} A genuine caller occasionally loses
consecutive windows on a live call that the same clip never loses offline. We tested
the obvious explanation --- that the browser's codec degrades the audio --- by
encoding each clip through that carrier offline and rescoring it, and it does not
hold: both genuine clips keep every window. By elimination the cause lies in the
real-time transport, and we have not isolated it.

\textbf{Interpretability covers two systems of three.} The probes could not be run on
the fully fine-tuned system, and the probe set is five clips: enough to refute an
explanation, not to characterise one.

\subsection{Future work}

The most valuable next result is the one that would change a deployment decision:
\emph{what makes an unseen tool detectable over a narrowband channel?} Our probing
suggests a direction --- adaptation moved the model's evidence for genuine speech
into the pitch-harmonic band that survives a codec --- so training explicitly for
band-limited evidence, rather than merely training on band-limited audio, is worth
testing. The augmentation lead of Sec.~\ref{sec:variants} should be confirmed on a
second unseen tool, and it carries a lesson for model selection: the in-domain dev set
did not register a 9-point gain on the unseen tool, so choosing among detectors needs
a validation split built from a synthesis tool held out from training, separate from
the one used for testing.

Beyond that: exercising the telephony carrier on live numbers; calibrating a second
operating point on real calls rather than on a simulated carrier; extending the
corpus beyond one code-mixed pair and one recording domain, since the present gap
measures \emph{English studio to Hindi--English lecture} and not code-mixing in
isolation; and harder attacks, since a detector is only ever as tested as the fakes
it was shown.
